\documentclass[11pt]{letter}
\usepackage[T1]{fontenc}
\usepackage[utf8]{inputenc}
\usepackage{lmodern}
\usepackage[a4paper,margin=25mm]{geometry}
\usepackage{url}
\signature{Ku Kang\\CBRN Defense Research Institute, Korea Ministry of National Defense\\
Seoul, Republic of Korea\\\texttt{bisu9082@gmail.com}}
\address{}
\begin{document}
\begin{letter}{The Editors\\Atmospheric Research}
\opening{Dear Editors,}

We submit the research paper \emph{``The CAMS reanalysis follows column dust in time but not
surface PM$_{10}$ across the Abu Dhabi monitoring network''} by Jin Yoo, Jeongyun Kim, Doo-Hee Lee
and Ku Kang for consideration in \emph{Atmospheric Research}.

Aerosol reanalyses assimilate optical depth, yet they are used to estimate surface coarse
particulate matter where monitors are few. We evaluate the CAMS reanalysis against the 19
regulatory PM$_{10}$ monitors of the Emirate of Abu Dhabi over 19 months, with an AERONET sun
photometer and airport weather reports as independent records of dust. Three results may interest
readers of the journal. First, the reanalysis times dust correctly in the column (its dust optical
depth follows the AERONET coarse-mode optical depth with Spearman 0.85, and its total optical depth
matches the measured one), but it gives dust too small a share of the column and carries too
little surface PM$_{10}$ per unit of column dust, so its regional deficit grows from 14 to
32~$\mu$g\,m$^{-3}$ with dust. Second, over the conurbation a local deficit of about
50--60~$\mu$g\,m$^{-3}$, present on clear and low-dust days and unrelated to modelled wind and dust,
adds to it; up to half of the between-monitor variance lies inside single model cells. Third, the
PM$_{10}$ values that an open archive holds for this network are 24-h running means; compared with
instantaneous model values they understate the correlation and produce an apparent benefit of
statistical correction that is not established when the two are averaged alike. The study extends
evaluations of reanalysis aerosol published in the journal, including that of its vertical
structure against spaceborne lidar (Zhang et al., 2025, Atmospheric Research 322, 108126), to a
dense surface network in the Arabian Gulf.

\textbf{Related manuscript.} A separate manuscript by the same authors is under consideration in
\emph{Atmospheric Environment} (ATMENV-D-26-01869R1, ``Meteorological covariates can mislead
zero-shot PM$_{2.5}$ nowcasting during Gulf dust extremes: an observation-validated benchmark'').
The two manuscripts share no data, model output or result. That study forecasts PM$_{2.5}$ from
the MERRA-2 aerosol reanalysis and ERA5 meteorology and validates it against OpenAQ monitors in
Manama, Kuwait City, Dhahran, Jeddah and Doha. The present study evaluates PM$_{10}$ from the
CAMS EAC4 reanalysis against the Abu Dhabi regulatory network, none of whose monitors appears in
the other manuscript, and involves no forecasting.

This manuscript has not been published and is not under consideration elsewhere. All authors
have approved the submission. All data, code and result files are openly available at
\url{https://github.com/bisu9082/camscast} and archived at Zenodo. The authors declare no
competing interests.

\closing{Sincerely,}
\end{letter}
\end{document}
